\documentclass[12pt,a4paper]{report}
\usepackage[utf8]{inputenc}
\usepackage[T1]{fontenc}
\usepackage{amsmath, amssymb, amsfonts}
\usepackage{geometry}
\usepackage{listings}
\usepackage{xcolor}
\usepackage{graphicx}
\usepackage{hyperref}
\usepackage{bm}
\usepackage[T1]{fontenc}
\usepackage[utf8]{inputenc}
\usepackage{listings}
\usepackage{titlesec}

\titleformat{\chapter}
  {\normalfont\huge\bfseries}
  {\thechapter.}
  {0.5em}
  {}
\titlespacing*{\chapter}
  {0pt}
  {-15pt}
  {20pt}

\lstset{
  inputencoding=utf8,
  extendedchars=true,
  literate=
    {á}{{\'a}}1 {à}{{\`a}}1 {ã}{{\~a}}1 {â}{{\^a}}1
    {é}{{\'e}}1 {ê}{{\^e}}1
    {í}{{\'i}}1
    {ó}{{\'o}}1 {õ}{{\~o}}1 {ô}{{\^o}}1
    {ú}{{\'u}}1
    {ç}{{\c{c}}}1
    {Á}{{\'A}}1 {É}{{\'E}}1 {Í}{{\'I}}1
    {Ó}{{\'O}}1 {Ú}{{\'U}}1
    {Ç}{{\c{C}}}1
}

\geometry{margin=2.5cm}

% Estilização de blocos de código
\definecolor{codegreen}{rgb}{0,0.6,0}
\definecolor{codegray}{rgb}{0.5,0.5,0.5}
\definecolor{codepurple}{rgb}{0.58,0,0.82}
\definecolor{backcolour}{rgb}{0.95,0.95,0.92}

\lstset{
    language=Octave,
    backgroundcolor=\color{backcolour},   
    commentstyle=\color{codegreen},
    keywordstyle=\color{magenta},
    numberstyle=\tiny\color{codegray},
    stringstyle=\color{codepurple},
    basicstyle=\ttfamily\small,
    breakatwhitespace=false,         
    breaklines=true,                 
    captionpos=b,                    
    keepspaces=true,                 
    numbers=left,                    
    numbersep=5pt,                  
    showspaces=false,                
    showstringspaces=false,
    showtabs=false,                  
    tabsize=2,
    frame=single
}

\title{Análise Matemática \\ \Large Projeto 3 - Simulação do Percurso Solar (Costa de Sagres)}
\author{
    \textbf{Grupo Nº: 2} \\    
    David Domingues (2022125309) \\
    Diana Gomes (2019101181) \\
    Guilherme Baltazar (2025209893) \\
    Joana Venâncio (2016123347) \\
    Leonardo Costa (2023137263) \\
    Marcelo Marques (2022102560) \\
    Miguel Reis (2024164633) \\ \\
    Escola Superior de Tecnologia e Gestão \\
    Instituto Politécnico de Leiria
}\date{Ano Letivo 2025/2026}

\begin{document}

\maketitle

\tableofcontents

\chapter{Fundamentação Teórica}

\section{Introdução à Geometria Solar}
A Terra é modelada como uma esfera unitária $S^2 = \{ \mathbf{x} \in \mathbb{R}^3 : \|\mathbf{x}\| = 1 \}$. Para um observador numa localização específica $\mathbf{p} \in S^2$ (ex: Sagres: $37.0097^\circ$N, $8.9408^\circ$W), a posição do Sol é definida pelo vetor unitário $\mathbf{s}(t)$.

\section{Projeção no Plano Tangente}
A direção observável do Sol a partir do ponto $\mathbf{p}$ é a projeção de $\mathbf{s}(t)$ no plano tangente a esse ponto, denotado como $T_p(S^2)$. 
A definição matemática do plano tangente é:
\begin{equation}
T_p(S^2) = \{ \mathbf{v} \in \mathbb{R}^3 : \mathbf{v} \cdot \mathbf{p} = 0 \}
\end{equation}
A projeção $\mathbf{v}_p(t)$ é obtida removendo a componente de $\mathbf{s}(t)$ que é normal à esfera em $\mathbf{p}$:
\begin{equation}
\mathbf{v}_p(t) = \mathbf{s}(t) - (\mathbf{s}(t) \cdot \mathbf{p}) \mathbf{p}
\end{equation}

\section{Representação em Toro Discretizado}
Para simular o percurso solar ao longo de um ano, mapeamos as coordenadas temporais (dia $d$ e minuto $m$) num toro $\mathbb{T}^2 = S^1 \times S^1$.
\begin{itemize}
    \item $d \in \{1, \dots, 365\}$ (Dias do ano)
    \item $m \in \{1, \dots, 1440\}$ (Minutos do dia)
\end{itemize}

\section{Níveis de Aproximação}
A precisão da simulação depende do modelo utilizado para a declinação solar ($\delta$) e para a Equação do Tempo ($E$):
\begin{enumerate}
    \item \textbf{Aproximação Grosseira:} Assume uma órbita circular e declinação constante. $E(d) = 0$.
    \item \textbf{Média (Linear):} Utiliza um modelo linear para $\delta$ e uma correção de primeira ordem para a Equação do Tempo.
    \item \textbf{Alta (Segunda Ordem):} Incorpora a excentricidade orbital e uma Equação do Tempo completa de segunda ordem para considerar a velocidade não uniforme da Terra em torno do Sol.
\end{enumerate}

\chapter{Exemplos Práticos e Implementação}

\section{Script de Simulação em Octave}
O seguinte código implementa a aproximação de alta ordem e visualiza os resultados numa malha toroidal.

\begin{lstlisting}[language=Octave, caption={Implementação em Octave para Simulação do Percurso Solar}]
% Coordenadas Geográficas de Sagres
lat = deg2rad(37.0097);
lon = deg2rad(-8.9408);

% Malha (Meshgrid) para o Toro (Dias vs Minutos)
d = 1:5:365;      % Amostragem a cada 5 dias para clareza
m = 0:30:1410;    % Amostragem a cada 30 minutos
[D, M] = meshgrid(d, m);

% 1. Aproximação Alta: Equação do Tempo (E em minutos)
B = (360/365) * (D - 81);
E = 9.87*sin(deg2rad(2*B)) - 7.53*cos(deg2rad(B)) - 1.5*sin(deg2rad(B));

% 2. Declinação (delta)
delta = deg2rad(23.45 * sin(deg2rad((360/365)*(284 + D))));

% 3. Tempo Solar e Angulo Horario (xomega)
solar_time = (M/60) + (E/60); 
omega = deg2rad(15 * (solar_time - 12));

% 4. Cálculo da Altitude (alpha)
sin_alpha = sin(lat)*sin(delta) + cos(lat)*cos(delta).*cos(omega);
alpha = asin(sin_alpha);
alpha(alpha < 0) = 0; % Noite: define altitude como 0

% 5. Visualização no Toro
R = 10; r = 4;
theta = (D/365) * 2 * pi;
phi = (M/1440) * 2 * pi;
X = (R + r*cos(phi)) .* cos(theta);
Y = (R + r*cos(phi)) .* sin(theta);
Z = r * sin(phi);

surf(X, Y, Z, rad2deg(alpha));
shading interp; colorbar;
title('Campo de Altitude Solar num Toro Discretizado');
\end{lstlisting}

\chapter{Exercícios Práticos}

\section{Exercício 1: Ortogonalidade Vetorial}
Prove analiticamente que o vetor $\mathbf{v}_p(t) = \mathbf{s}(t) - (\mathbf{s}(t) \cdot \mathbf{p})\mathbf{p}$ é sempre ortogonal a $\mathbf{p}$. 
\textit{Sugestão: Calcule o produto escalar $\mathbf{v}_p(t) \cdot \mathbf{p}$ e utilize o facto de que $\|\mathbf{p}\|^2 = 1$.}

\subsection{Resolução}
Queremos provar que
\[
\mathbf{v}_p(t)=\mathbf{s}(t)-(\mathbf{s}(t)\cdot \mathbf{p})\mathbf{p}
\]
é ortogonal a \(\mathbf{p}\).
Para isso, calculamos o produto escalar:
\[
\mathbf{v}_p(t)\cdot \mathbf{p}
=\left[\mathbf{s}(t)-(\mathbf{s}(t)\cdot \mathbf{p})\mathbf{p}\right]\cdot \mathbf{p}.
\]
Usando a distributividade do produto escalar, obtemos:
\[
\mathbf{v}_p(t)\cdot \mathbf{p}
=\mathbf{s}(t)\cdot \mathbf{p}
-(\mathbf{s}(t)\cdot \mathbf{p})(\mathbf{p}\cdot \mathbf{p}).
\]
Como \(\mathbf{p}\in S^2\), então \(\|\mathbf{p}\|=1\), logo:
\[
\mathbf{p}\cdot \mathbf{p}=\|\mathbf{p}\|^2=1.
\]
Substituindo:
\[
\begin{aligned}
\mathbf{v}_p(t)\cdot \mathbf{p}
&=\mathbf{s}(t)\cdot \mathbf{p}
-(\mathbf{s}(t)\cdot \mathbf{p})\cdot 1.
\end{aligned}
\]
Portanto:
\[
\mathbf{v}_p(t)\cdot \mathbf{p}=0.
\]
Concluímos que:
\[
\mathbf{v}_p(t)\perp \mathbf{p}.
\]

Assim, o vetor \(\mathbf{v}_p(t)\) pertence ao plano tangente \(T_p(S^2)\).
\newpage
\section{Exercício 2: Análise de Sensibilidade}
Compare a altitude solar $\alpha$ calculada para o solstício de verão ($d=172$) às 12:00h utilizando a aproximação "Grosseira" ($E=0$) e a aproximação "Alta". Qual é o erro em graus?

\subsection{Resolução}

Pretende-se comparar a altitude solar \(\alpha\) no solstício de verão, correspondente a \(d=172\), às \(12:00\), usando duas aproximações.
A latitude de Sagres é:
\[
\varphi = 37.0097^\circ.
\]
A altitude solar é dada por:
\[
\sin(\alpha)
=
\sin(\varphi)\sin(\delta)
+
\cos(\varphi)\cos(\delta)\cos(\omega),
\]
onde \(\delta\) é a declinação solar e \(\omega\) é o ângulo horário.
No solstício de verão, toma-se aproximadamente:
\[
\delta \approx 23.45^\circ.
\]
\subsection*{Aproximação Grosseira}
Na aproximação grosseira, considera-se:
\[
E(d)=0.
\]
Às \(12:00\), o tempo solar é:
\[
t_s=12.
\]
Logo:
\[
\omega = 15^\circ(t_s-12)=0^\circ.
\]
Assim:
\[
\begin{aligned}
\sin(\alpha_g)
&=
\sin(37.0097^\circ)\sin(23.45^\circ) \\
&\quad+
\cos(37.0097^\circ)\cos(23.45^\circ)\cos(0^\circ).
\end{aligned}
\]
Como \(\cos(0^\circ)=1\), temos:
\[
\alpha_g \approx 76.44^\circ.
\]
\subsection*{Aproximação Alta}
Na aproximação alta, usa-se a Equação do Tempo:
\[
B=\frac{360}{365}(d-81).
\]
Para \(d=172\):
\[
B=\frac{360}{365}(172-81)\approx 89.75^\circ.
\]
A Equação do Tempo é:
\[
E = 9.87\sin(2B)-7.53\cos(B)-1.5\sin(B).
\]
Substituindo:
\[
\begin{aligned}
E
&= 9.87\sin(2B)-7.53\cos(B)-1.5\sin(B) \\
&\approx -1.39 \text{ minutos}.
\end{aligned}
\]
O tempo solar corrigido é:
\[
\begin{aligned}
t_s
&= 12+\frac{E}{60} \\
&\approx 12-\frac{1.39}{60} \\
&\approx 11.9768.
\end{aligned}
\]
Logo:
\[
\begin{aligned}
\omega
&= 15^\circ(t_s-12) \\
&\approx 15^\circ(-0.0232) \\
&\approx -0.348^\circ.
\end{aligned}
\]
A altitude solar é então:
\[
\begin{aligned}
\sin(\alpha_a)
&=
\sin(37.0097^\circ)\sin(23.45^\circ) \\
&\quad+
\cos(37.0097^\circ)\cos(23.45^\circ)\cos(-0.348^\circ).
\end{aligned}
\]
Daqui resulta:
\[
\alpha_a \approx 76.44^\circ.
\]
\subsection*{Erro entre as aproximações}
O erro absoluto é:
\[
\text{Erro}=|\alpha_g-\alpha_a|.
\]
Substituindo:
\[
\text{Erro}\approx |76.44^\circ-76.44^\circ|.
\]
Numericamente:
\[
\text{Erro}\approx 0.0007^\circ.
\]
Portanto, para o solstício de verão às \(12:00\), a diferença entre a aproximação grosseira e a aproximação alta é praticamente desprezável.
\newpage
\section{Exercício 3: Integração no Toro}
Dada a função de altitude $\alpha(d, m)$, monte o integral duplo sobre o toro para encontrar a altitude solar média durante as horas de sol para o mês de julho.

\subsection{Resolução}
Pretende-se montar o integral duplo que permite calcular a altitude solar média durante as horas de sol no mês de julho.
A função de altitude solar é:
\[
\alpha=\alpha(d,m),
\]
onde:
\[
d
\]
representa o dia do ano e
\[
m
\]
representa o minuto do dia.
O mês de julho corresponde aproximadamente aos dias:
\[
d\in[182,212].
\]
As horas de sol são aquelas em que a altitude solar é positiva:
\[
\alpha(d,m)>0.
\]
Assim, define-se o domínio de integração:
\[
\begin{aligned}
D_{\text{sol}}
&=
\{(d,m)\in[182,212]\times[0,1440] : \alpha(d,m)>0\}.
\end{aligned}
\]
A altitude solar média durante as horas de sol é dada por:
\[
\overline{\alpha}
=
\frac{1}{A(D_{\text{sol}})}
\iint_{D_{\text{sol}}} \alpha(d,m)\,dm\,dd,
\]
onde \(A(D_{\text{sol}})\) representa a área do domínio correspondente às horas de sol:
\[
A(D_{\text{sol}})
=
\iint_{D_{\text{sol}}} 1\,dm\,dd.
\]
Portanto:
\[
\boxed{
\overline{\alpha}
=
\frac{
\displaystyle \iint_{D_{\text{sol}}} \alpha(d,m)\,dm\,dd
}{
\displaystyle \iint_{D_{\text{sol}}} 1\,dm\,dd
}
}
\]
Como a variável temporal é representada num toro \(\mathbb{T}^2\), podemos também escrever o integral usando as coordenadas angulares:
\[
\theta=\frac{2\pi d}{365},
\qquad
\phi=\frac{2\pi m}{1440}.
\]
Nesse caso, o domínio correspondente a julho é:
\[
\theta\in
\left[
\frac{2\pi\cdot 182}{365},
\frac{2\pi\cdot 212}{365}
\right].
\]
A altitude média pode ser escrita como:
\[
\overline{\alpha}
=
\frac{
\displaystyle \iint_{\widetilde{D}_{\text{sol}}}
\alpha(\theta,\phi)\,d\phi\,d\theta
}{
\displaystyle \iint_{\widetilde{D}_{\text{sol}}}
1\,d\phi\,d\theta
}.
\]
onde:
\[
\begin{aligned}
\widetilde{D}_{\text{sol}}
&=
\{(\theta,\phi): \alpha(\theta,\phi)>0\}.
\end{aligned}
\]
Este integral representa a média da altitude solar apenas nos pontos do toro correspondentes ao mês de julho e aos momentos em que o Sol está acima do horizonte.

\chapter{Problemas Sugeridos para a Defesa Oral}

As seguintes questões foram desenhadas para avaliar o domínio do projeto pelo grupo:

\begin{enumerate}
    \item \textbf{Interpretação Geométrica:} Explique como o mapeamento $(d, m) \to \mathbb{T}^2$ preserva a continuidade do tempo (ou seja, como a transição de 31 de dez para 1 de jan é gerida pela topologia do toro).
    \item \textbf{Aplicação de Cálculo:} Como usaria o \textit{Gradiente} da função de altitude $\alpha(d, m)$ para encontrar o momento exato do ano em que o sol nasce mais rapidamente em Sagres?
    \item \textbf{Métodos Numéricos:} No código Octave, utilizámos um \textit{meshgrid}. Discuta como o tamanho do passo de $d$ e $m$ afeta a precisão da simulação e o custo computacional.
    \item \textbf{Transformação de Coordenadas:} O vetor $\mathbf{s}(t)$ é frequentemente dado em coordenadas equatoriais. Explique as matrizes de rotação necessárias para o transformar no sistema de coordenadas horizontais local no ponto $\mathbf{p}$.
    \item \textbf{Pensamento Crítico:} Se a órbita da Terra fosse perfeitamente circular, qual nível de aproximação seria suficiente para modelar o percurso do Sol? Justifique com base na Equação do Tempo.
    \item \textbf{Colaboração com IA:} Identifique uma parte do código Octave fornecido pela IA que exigiu correção manual ou verificação matemática pelo grupo para garantir que seguia as restrições específicas do projeto.
\end{enumerate}

\end{document}